% !TeX root = ../papers/nc-p3.tex
\section{中心の滑らかな部分と二十個の最小単純指標}
\label{sec:minimum-central-characters}
ここでは前二節の正規中心とPI次数16を持つ開集合を扱う。
体は代数閉、標数零とする。$B$は折り畳み環、$Z$は全中心、
$P$は四つの長さ4の中心変数の生成する環である。

\begin{proposition}\label{prop:smooth-azumaya}
$B$のAzumaya軌跡は$\operatorname{Spec}Z$の滑らかな部分に一致する。
$(u,v)\mapsto\operatorname{trd}(uv)$は$B\to\operatorname{Hom}_Z(B,Z)$
の同型を誘導する。ここで$\operatorname{trd}$は中心単純な一般商の被約traceである。
\end{proposition}
\begin{proof}
$B$は$P$上自由で、$\bar B=B/P_+B$は最高経路長12の斉次Frobenius代数である。
後者を確認する。正則中心列四元を一つずつ除くと、各頂点単純$S_i$に対し
$\operatorname{Hom}_{\bar B}(S_i,\bar B)$は内部次数の移動を除き
$\operatorname{Ext}^4_B(S_i,B)$に一致する。一元の場合は、非零因子の乗法による
二項分解から得る環変更のExtが一行に集中する。四段では変数の次数和16と
AS指数4の差により、各頂点に対応する左socleは経路長12に一つずつある。
このsocle上で各頂点成分に非零な長さ12の線型汎関数を選ぶと、その核は
非零左イデアルを含まない。実際、各単純socleは相異なる頂点型で重複しない。
従って積と汎関数から得る対合は完全である。

この汎関数を$P$自由な斉次基底で、経路長を12だけ下げる$P$線型汎関数へ
持ち上げる。積対合で定まる左$B$加群写像は斉次であり、$P_+$剰余で同型なので
次数付き中山と自由階数の一致から同型となる。一方$Z$は$P$上の四自由基底の
最後が長さ12で、その係数を取る写像が$Z\simeq\operatorname{Hom}_P(Z,P)$を
同じ次数移動付きで与える。実際、$P_+$剰余は四次モニック一元商または
重み付き二次・四次完全交叉で、最高次のsocleが一次元である。
随伴と二つの移動の相殺から、$\operatorname{Hom}_Z(B,Z)$と$B$は
次数移動なしに左$B$加群として同型である。

被約traceは$Z$に値を取る。各元は$Z$上整であり、分裂体での固有値の和は整、
かつ中心の分数体に入り、$Z$は正規だからである。次数作用との可換性により
trace対合の写像は次数0で、一般商で非退化である。先の左加群同型と合成すると
$B$の次数0の左加群自己準同型、すなわち$\sum_i c_i e_i$の右乗法になる。
一般商で可逆なので全$c_i\ne0$であり、元の写像も同型である。

$B$は有限な$Z$加群として極大Cohen--Macaulayである。
各極大点ではその下の$P$の四局所パラメータが正則列であり、局所化でもこの性質を保つ。
従って正則な中心局所環上ではAuslander--Buchsbaum公式により自由となる。
自由基底と被約trace双対基底を$u_j,v_j$とする。一般商での行列単位の計算から
$\sum_j u_j\otimes v_j$は両側作用を交換し、積の像は$16\cdot1$である。
両側の自由性によりこの等式は局所環上で成立する。16で割った元が分離元を与え、
中心が基底環に等しい有限射影的分離代数はAzumayaとなる。
逆にAzumaya点では$B$の局所化は中心上自由である。Koszul両側分解から
$B$の通常の大域次元は高々4であり、中央局所化後もこの上界を持つ。
ここで両側分解は$B\otimes_{B_0}W_j\otimes_{B_0}B$に左右の端文字の微分を置いたもの
で、各全次数で右因子次数により濾過すると片側Koszul完全性から対角$B$を解消する。
局所中心の剰余体を$k'$とし、$B\otimes_Z k'$の射影分解を中心へ制限する。
$k'$はこの自由な$k'$空間の直和因子なので中心上有限射影次元を持ち、
局所中心は正則である。
\end{proof}

\begin{theorem}\label{thm:twenty-minimal-characters}
中心錐の頂点外にある各単純右加群の頂点次元は$(d,d,d,d)$で、
$d\in\{1,2,3,4\}$である。十五次元成分の一般窓では、$d=1$の同型類は
次数によるスカラー作用を除くと丁度二十個であり、対応する射影中心上の点は
全て異なる特異点である。各点の一つの非零アフィン代表上には丁度二種類の
単純加群があり、次元4のものと、次元8または12のものからなる。
\end{theorem}
\begin{proof}
任意の頂点外の中心指標ではある$x\in P_4$が可逆である（添字4は経路長）。
実際$Z$は$P$上斉次有限である。$S$をこの指標の単純右加群とし、
$M_n=Se_{n\bmod4}$を$n\ge0$に置き、矢を$S$の作用で掛ける。
$x$は$M_n\to M_{n+4}$の同型を与え、$M$は次数付き有限生成である。
$B$はNoetherで、有限長の次数付き射影分解の各項は頂点射影の移動の有限直和である。
各頂点射影の全Hilbert級数は$(1-t)^{-4}$なので
\[
 H_M(t)=\frac{d_0+d_1t+d_2t^2+d_3t^3}{1-t^4}
       =\frac{h(t)}{(1-t)^4},\qquad d_i=\dim Se_i .
\]
非自明な四乗根三個で分子が消え、四次元数は等しい。PI次数16または
一般行列環の標準恒等式と単純像の全行列性から$4d\le16$である。

$d=1$なら各辺の四矢の値は同時に零ではない。逆に、四辺で非零な一次元空間を
結ぶ矢組は有向閉路をなすため単純である。関係式は各辺の双一次式そのものである。
第一辺は有限エタールな二十点、残りの三辺はグラフで、四移動は恒等である。
従って射影化した四組は丁度二十個で、スカラーを戻すと各組につき四辺の倍率は
頂点基底変換によりその積一つを残す。次数作用でこの積を任意の非零値へ移せる。
各単純は頂点外である。頂点指標の商は有限次元の正次数付き環であり、
正次数イデアルが冪零だから単純は四つの頂点単純に限られる。
一つの非零な$x$の値を固定すれば積が決まり、一つの射影組から一種類となる。

二十指標の分離を検算する。$t$の二つの切断の法線写像の共通核から、
節\ref{sec:unrestricted-doublet}の有理商接代表十五個の固定整数線型結合$T$を選ぶ。
$p=2657$では第一辺が完全分裂する。$W+pT$から二次関係を取り、法$p^2$の
経路長5までの商基底で四次の可換核の自由基底五元を求める。
各射影四組では各位置で一座標を1とし、残り十二座標の修正を二十四本の
辺方程式の微分で解く。全てで十二ピボットがあり、修正後の全式が消える。
第0隅の中心五元の各経路を四矢の積で評価し、単元座標で正規化すると、
剰余体では八値、法$p^2$では二十値が異なる。
\path{computations/verify_minimum_central_characters.py}が同じ接線と番号付けで再現する。

この分離は標数零の一般窓に移る。番号付き辺はエタールで、$t$の成分は滑らかなので、
標数零の接ベクトル$T$はこの被覆へ一意に持ち上がる。分裂素数での各番号は
$\mathbb Z_p$上へHenselで延びる。中心と経路の基底は検算の単元ピボットで取り、
一次微分の持ち上げも$\mathbb Z_p[\epsilon]/(\epsilon^2)$上で計算できる。
実際、余分な行の等式は$\mathbb Q_p[\epsilon]/(\epsilon^2)$上で
幾何学的な一次変形により成立し、単元による解の係数は整である。
$\epsilon\mapsto p\bmod p^2$による像が上の検算である。
よってどの二切断も近傍で恒等的には一致しない。番号付けは有限エタールであり、
対角を除いた対の指標一致軌跡の像を除けば一般窓で二十点全てを分離できる。
次数4の単純はAzumaya点の単純次元16より小さいため、各中心指標は前命題により
特異で、その射影点も特異となる（$x$を可逆にした中心は換算次数1の可逆元を持つ）。

最後に、一つの射影点上で$x$の値を1に規格化する。最小単純$S$の自己一次変形は
$B/(x-1)$上では全て自明である。実際、頂点基底を選んだ双対数上の変形でも
各頂点は自由階数1で、二十点のエタール性とグラフ性から各射影矢は固定である。
四倍率の積は$x=1$により固定で、残りは基底変換で消せる。
この三次元の中心を指標で完備化しても同じ自己拡大の消滅が成立する。
長さ2の拡大は中心極大イデアルの二乗で消え、完備化前後で同じだからである。
完備環上有限な代数は半完全である。もし単純が$S$しかなければ、その森田同値な
基本代数の根基$J$は$J/J^2=0$となる（これは単一単純の自己拡大の双対）。
中山により$J=0$で、中心が正次元であることに矛盾する。
ここで$x=1$近傍の中心は三次元の射影チャートである。例えば中心の$x$を
可逆にしてその四乗根をエタールに加えると、経路次数の作用で
$B/(x-1)$と可逆な一変数のテンソル積となる。中心、一般商の次数、有限性は
この自明化と両立する。

この完備中心は整域である（標数零の正規な多様体の局所環は優秀で正規性が完備化で保たれる）。
完備代数はこの環上torsion-freeで、一般商は次数16の中心単純環である。
各単純の次元を$n_j$とすると、半単純剰余の各行列環の原始対角冪等元は
直交した完備系へ持ち上がり、同じ行列環内では同値に持ち上がる。
一般分裂体で各原始元は正の行列階数$m_j$を持つので
$16=\sum_j n_jm_j$である。$n_j=4d_j$で割り直すと$4=\sum_j d_jm_j$。
$d=1$の単純は指標につき一つであり、ほかに少なくとも一つあるから、
残る可能性は$(d,m)=(1,1),(3,1)$または$(1,2),(2,1)$だけである。
\end{proof}
